\newpage

\section{\huge \textcolor{blue}{\hspace*{0.2em}线性非齐次微分方程组}}
\subsection{\LARGE \hspace*{0.1em}常系数非齐次微分方程组}
\subsubsection{\large 定理及性质}
标准形式:
\[\vec{X'} =   A \vec{X} \]
定理A:\\
上标准方程组的通解为:
\[\vec{x} = c_1 \vec{x}_1 + c_2 \vec{x}_2  \]
其中: $\vec{x}_1 ,\vec{x}_2$是两个线性无关解.\\[1em]
定理B:\\
朗斯基行列式
\[ w(\vec{x}_1 ,\vec{x}_2) = 
\begin{vmatrix}
    \vec{x}_1 \vec{x}_2
\end{vmatrix}
\]
\[ w(t) = 
\begin{cases}
    \equiv 0 \hspace*{1em} \text{if $\vec{x}_1 ,\vec{x}_2$ 相关} \\
    \ne 0 \hspace*{1em} \text{if $\vec{x}_1 ,\vec{x}_2$ 无关}
\end{cases}
\]
同时所对应的矩阵也称为方程组的基本矩阵
\[\vec{X'} =   A \vec{X} \]
\[ \bar{X} =
\begin{bmatrix}
    \vec{x}_1  \vec{x}_2
\end{bmatrix}
\hspace*{1em} \vec{x}_1 ,\vec{x}_2 \text{ 无关解}
\]
性质1: $\left\lvert\bar{X}\right\rvert  \ne 0$ \\
性质2:
\begin{align*}
\vec{X'} = A \vec{X} &\Leftrightarrow \left[ \vec{x}_1'  \vec{x}_2'\right]   \\
    &= A \left[ \vec{x}_1  \vec{x}_2\right]  \\
    &= \left[ A \vec{x}_1  A\vec{x}_2\right] \\
    &\Leftrightarrow   \vec{x}_1'= A \vec{x}_1 , \vec{x}_2'= A \vec{x}_2
\end{align*}
现将上述定理及性质应用于解非齐次方程组上
\[
\begin{cases}
    x' = ax + by + r_1(t) \\
    y' = cx + dy + r_2(t)
\end{cases}
\]
\[\vec{X'} = A \vec{X} + \vec{r}(t) \]
定理C: $\vec{X'}$通解
\[ \vec{X} = \vec{X}_c + \vec{X}_p \]
其中: $\vec{X}_c$为对应齐次方程组通解, $\vec{X}_p$为方程组特解.\\
$\vec{X}_c$在前文已作相应讨论, 下重点讨论$\vec{X}_p$\\[1em]
利用参数变分法:
\begin{align*}
    \vec{X}_p &= v_1(t) \vec{x}_1 + v_2(t) \vec{x}_2 \\
    \vec{X}_p &= \bar{X}  \vec{V} \\
    \vec{X}_p' &= A \vec{X}_p + \vec{r} 
\end{align*}
\[
    \bar{X}'\vec{V} + \bar{X}\vec{V}'=A \vec{X}\vec{V} + \vec{r}
\]
$\because \bar{X}'= A \vec{X}$
\begin{align*}
    \bar{X}\vec{V}' &= \vec{r} \\
    \vec{V}' &= \bar{X}^{-1} \vec{r}
\end{align*}
$\bar{X}$的逆矩阵存在, 需行列式$\neq 0$
\begin{align*}
    \vec{V} &=  \int \bar{X}^{-1} \vec{r} dt \\
    \vec{X}_p &= \bar{X} \int \bar{X}^{-1} \vec{r} dt
\end{align*}

\newpage
\subsection{\LARGE \hspace*{0.1em}矩阵指数及应用}
\subsubsection{\large 定理及性质}
上文提到A的基本矩阵性质\\
现考虑: 基本矩阵不唯一, 解的选取仅需线性无关\\
方程通解:
\[ \vec{x} = c_1 \vec{x}_1 + c_2 \vec{x}_2 \]
\[ \vec{X} = \mathbb{X} \vec{C} =
\begin{bmatrix}
     \vec{x}_1  \vec{x}_2
\end{bmatrix}\cdot 
\begin{bmatrix}
    c_1 \\
    c_2
\end{bmatrix}
\]
\[
\begin{bmatrix}
    \mathbb{X} \vec{c}_1  & \mathbb{X}\vec{c}_2
\end{bmatrix}
=
\mathbb{X}
\begin{bmatrix}
    \vec{c}_1  & \vec{c}_2
\end{bmatrix}
=
\mathbb{X} \vec{C} \hspace*{1em} \left\lvert C\right\rvert \neq 0
\]
通解公式:\\
对一阶微分方程:
\[ x' = a x \hspace*{1em} , x = c e^{at} \]
即猜测在解方程组通解也会利用指数函数性质, 为此先对矩阵指数定义:\\
借用无穷级数的指数函数定义:
\begin{align*}
e^{at} &= 1 + at + \frac{a^2 t^2}{2!} + \frac{a^3 t^3}{3!} + \cdots    \\
\frac{de^{at}}{dt} &= a + a^2 t + \frac{a^3}{2!} t^2 + \cdots \\
    &=a (1 + at + \frac{a^2 t^2}{2!} + \frac{a^3 t^3}{3!} + \cdots)  \\
    &=a e^{at}
\end{align*}
$\Downarrow $猜测
\[ \vec{X}' = A \vec{X} \hspace*{1em}  \vec{X} = e^{At} \]
矩阵指数的定义:
\[ e^{At} = I + At + \frac{A^2 t^2}{2!} + \frac{A^3 t^3}{3!} + \cdots \]
验证是否为基本矩阵:\\
代入方程$\vec{X}' = A \vec{X} $, 显然成立\\
同时满足$\left\lvert X(0)\right\rvert  = \left\lvert I\right\rvert = 1 \neq 0$\\
而矩阵指数并不好计算, 下举一个能计算的例子:\\[1em]
eg:
\[
\begin{cases}
    x' = y \\
    y' = x
\end{cases}
\hspace*{1em}
\vec{X}' =
\begin{bmatrix}
    0 & 1 \\
    1 & 0
\end{bmatrix}
\hspace*{1em}
A =
\begin{bmatrix}
    0 & 1 \\
    1 & 0
\end{bmatrix}
\]
\begin{align*}
A &=
\begin{bmatrix}
    0 & 1 \\
    1 & 0
\end{bmatrix}\\
A^2 &=
\begin{bmatrix}
    1 & 0 \\
    0 & 1
\end{bmatrix} = I \\ 
A^3 &= A^2 A = I A = A \\
A^4 &= A^3 A = A A = I \\
\cdots \cdots
\end{align*}
即
\begin{align*}
    e^{At} &= I + At + \frac{A^2 t^2}{2!} + \frac{A^3 t^3}{3!} + \cdots \\
    &= I + 
\begin{bmatrix}
    0 & 1 \\
    1 & 0
\end{bmatrix}t +I \frac{t^2}{2!}  + 
\begin{bmatrix}
    0 & 1 \\
    1 & 0
\end{bmatrix}\frac{t^3}{3!}  + \cdots \\
&=
\begin{bmatrix}
    1+\frac{t^2}{2!} + \cdots & t + \frac{t^3}{3!} + \cdots \\
    t + \frac{t^3}{3!} + \cdots & 1+\frac{t^2}{2!} + \cdots
\end{bmatrix}\\
&=
\begin{bmatrix}
    \cosh t & \sinh t \\
    \sinh t & \cosh t
\end{bmatrix}
\end{align*}
即求解.
\newpage
现探求矩阵指数对解初值问题的便利:\\
初值问题?
\[
\begin{cases}
    \vec{X}' = A \vec{X} \\
    \vec{X}(0) = \vec{X}_0
\end{cases}
\Longrightarrow \vec{X}(t)
\]
通解:
\begin{align*}
    \vec{X} &= e^{At} \vec{C} \\
    \vec{X}(0) &= e^{A0} \vec{C} \\
    &= I \vec{C} \\
    &= \vec{X}_0 \\
    \therefore \hspace*{0.5em} \vec{X} &= e^{At} \vec{X}_0
\end{align*}
\textcolor{red}{原则性警告}
\[ e^{A+B} \neq e^{A} e^{B} \]
若成立, 当且仅当 $AB = BA$ (可交换矩阵) \\
满足下3种情况之一即为可交换矩阵:
\begin{enumerate}
    \item $A = C I =
\begin{bmatrix}
    c & 0 \\
    0 & c
\end{bmatrix}$
    \item $B=-A$
    \item $B=A^{-1}$
\end{enumerate}
对$B=-A$:
\begin{align*}
    e^{A-A} &= e^A e^{-A}\\
    I &= e^A e^{-A} \\
    (e^A)^{-1} &= e^{-A}
\end{align*}
\newpage
计算$e^{At}$的各方法:
\begin{enumerate}
    \item 级数(限制太大)
    \item 对称矩阵
    \item $\mathbb{X} \cdot \mathbb{X}(0)^{-1} = e^{At}$
\end{enumerate}
其中关于对称矩阵采用分解处理:
\[
\begin{bmatrix}
    a & b \\
    b & a
\end{bmatrix} = 
\begin{bmatrix}
    a & 0 \\
    0 & a
\end{bmatrix} +
\begin{bmatrix}
    0 & b \\
    b & 0
\end{bmatrix}
\]


\subsection{\LARGE \hspace*{0.1em}方程组解耦}
\subsubsection{\large 定理及性质}
针对$n\times n$齐次方程组, 有
\[\vec{X}' = A \vec{X} \]
解法:
\begin{enumerate}
    \item 特征值 , 特征向量
    \item $e^{At}\vec{X_0}$ , $e^{At}=\mathbb{X} \cdot \mathbb{X}(0)^{-1}$
    \item 方程组解耦, 即本节内容
\end{enumerate}
eg:
\[
\begin{cases}
    u=ax + by \\
    v=cx + dy
\end{cases}
\hspace*{1em} \Rightarrow u,v
\]
若得:
\[
\begin{cases}
    u' = \lambda_1 u \\
    v' = \lambda_2 v
\end{cases}
\]
即称为方程组解耦.\\
耦合意味着,方程组右侧都出现x,y;而此时可分别解$u,v$ 即方程组被解耦.\\
eg:
\[
\begin{cases}
    1\cdot x' = c (y - x) \\
    2\cdot y' = c (x - y)
\end{cases}
\]
为演示方法,现令$c=2$, 得:
\[
\begin{cases}
    x' = -2x + 2y  \\
    y' = x - y
\end{cases}
\]
令
\[
\begin{cases}
    u = x + 2y \\
    v = x - y
\end{cases}
\]
代入,得
\[
\begin{cases}
    u' = x' + 2y' = 0\\
    v' = x' - y' = -3v
\end{cases}
\]
\[
\therefore u = c_1 , v = c_2 e^{-3t}
\]
解:
\[
\begin{cases}
    x = \frac{1}{3}(u + 2v) = \frac{1}{3}(c_1 + 2c_2 e^{-3t}) \\
    y = \frac{1}{3}(u - v) = \frac{1}{3}(c_1 - c_2 e^{-3t})
\end{cases}
\]
\[
\vec{X} = \frac{1}{3} c_1
\begin{bmatrix}
    1 \\
    1
\end{bmatrix} +
\frac{1}{3} c_2
\begin{bmatrix}
    2 \\
    -1
\end{bmatrix}
e^{-3t}
\]
现进行一般化,需先明确解耦条件:\\
特征值必须都是实数且完备的.\\
eg:
\[
\begin{bmatrix}
    u \\
    v
\end{bmatrix} =
\begin{bmatrix}
    a_1' & a_2' \\
    b_1' & b_2'
\end{bmatrix}
\begin{bmatrix}
    x \\
    y
\end{bmatrix}
\]
\[
\begin{bmatrix}
    x \\
    y
\end{bmatrix}=
\begin{bmatrix}
    a_1 & a_2 \\
    b_1 & b_2
\end{bmatrix}
\begin{bmatrix}
    u \\
    v
\end{bmatrix}
\]
记$
\begin{bmatrix}
    a_1' & a_2' \\
    b_1' & b_2'
\end{bmatrix}$ 为$D$,解耦矩阵; 
$\begin{bmatrix}
    a_1 & a_2 \\
    b_1 & b_2
\end{bmatrix}$ 为$D^{-1}=E$\\
\[ E =
\begin{bmatrix}
    \vec{\alpha}_1 & \vec{\alpha}_2
\end{bmatrix}
\]
两列是两个特征向量\\
代入$\vec{X}' = A \vec{X}$
\[ A \vec{\alpha}_1 = \lambda_1 \vec{\alpha}_1 \]
\begin{align*}
    A E &= A [\vec{\alpha}_1 \hspace*{0.5em} \vec{\alpha}_2] \\
    &= [A \vec{\alpha}_1 \hspace*{0.5em} A \vec{\alpha}_2] \\
    &= [\lambda_1 \vec{\alpha}_1 \hspace*{0.5em} \lambda_2\vec{\alpha}_2] \\
    &= [\vec{\alpha}_1 \hspace*{0.5em} \vec{\alpha}_2] 
    \begin{bmatrix}
        \lambda_1 & 0 \\
        0 & \lambda_2
    \end{bmatrix} \\
    &= E
    \begin{bmatrix}
        \lambda_1 & 0 \\
        0 & \lambda_2
    \end{bmatrix}
\end{align*}
代替:
\[ \vec{X}' = A \vec{X} \hspace*{1em} \vec{U}=
\begin{bmatrix}
    u \\
    v
\end{bmatrix} \]
\[ \vec{X} = E \vec{U} \]
\begin{align*}
    E \vec{U}' &= A E \vec{U} \\
    &= E
    \begin{bmatrix}
        \lambda_1 & 0 \\
        0 & \lambda_2
    \end{bmatrix} \vec{U} \\
    \therefore 
    \vec{U}' &= 
    \begin{bmatrix}
        \lambda_1 & 0 \\
        0 & \lambda_2
    \end{bmatrix} \vec{U}
\end{align*}
解耦得;
\[
\begin{cases}
    u' = \lambda_1 u \\
    v' = \lambda_2 v
\end{cases}\hspace*{1em}
 \therefore
\begin{cases}
    u = c_1 e^{\lambda_1 t} \\
    v = c_2 e^{\lambda_2 t}
\end{cases}
\]

\newpage
\subsection{\LARGE \hspace*{0.1em}非线性自治方程组及作图}
\subsubsection{\large 定理及性质}

现有非线性方程组:
\[
\begin{cases}
    x' = f(x,y) \\
    y' = g(x,y)
\end{cases}
\]
如何画轨迹?\\
eg:非线性摆:\\
\begin{figure}[htbp]
    \centering
\includegraphics[width=0.5\textwidth]{C:/Users/37/Desktop/Notes on Common Calculus Equations/Notes and pictures/cwf12.4.1.png}
    \caption{非线性摆}
    \label{fig:cwf12.4.1}
\end{figure}

\begin{align*}
    m\vec{a} & = \vec{F} \\
    m l \theta '' & = - mg\sin\theta - c_1 l \theta' \\
    \theta '' + \frac{c_1}{m}\theta' &+ \frac{g}{l}\sin\theta  = 0 \\
    \theta '' + c\theta' &+ k\sin\theta  = 0 \\
    \theta ' &= \omega \\
    \omega ' &= - k\sin\theta - c\omega 
\end{align*}
令$c=1, k=2$, 得
\[
\begin{cases}
    \theta' = \omega \\
    \omega' = - 2\sin\theta -\omega  
\end{cases}
\]
1.观察方程组的临界点$(x_0,y_0)$
\[
\begin{cases}
    f(x_0,y_0)=0 \\
    g(x_0,y_0)=0
\end{cases}
\]
联立方程:
\[
\begin{cases}
    \omega = 0 \\
    - 2\sin\theta -\omega = 0 
\end{cases} \Rightarrow 
\begin{cases}
    \theta = k\pi \\
    \omega = 0
\end{cases} \hspace*{0.5em} k \in \mathbb{Z} 
\]
而由物理现象实际作限制 $\theta = 0 , \pi$\\
临界点:$(0,0)$ 稳定点 ; $(\pi,0)$ 不稳定点\\[1em]
2.对每个临界点附近线性化方程组\\
在$(0,0)$附近线性化方程:
\[
\begin{cases}
    \theta' = \omega \\
    \omega' = - 2\theta -\omega  
\end{cases}
\hspace*{0.5em}
\begin{bmatrix}
    0 & 1 \\
    -2 & -1
\end{bmatrix}
\]
\[ \lambda^2 + \lambda + 2 = 0 \]
\[ \lambda = \frac{-1 \pm \sqrt{-7}}{2} \]
因为$\lambda=-\frac{1}{2} + b i$,$Re < 0$, 图像将呈螺旋汇聚\\
\begin{figure}[htbp]
    \centering
\includegraphics[width=0.5\textwidth]{C:/Users/37/Desktop/Notes on Common Calculus Equations/Notes and pictures/cwf12.4.2.png}
    \caption{螺旋汇聚}
    \label{fig:cwf12.4.2}
\end{figure}

在$(\pi,0)$附近线性化方程:
\[
\mathcal{J} = 
\begin{bmatrix}
    f_x & f_y \\
    g_x & g_y
\end{bmatrix}_0
\]
利用雅可比矩阵在$(x_0,y_0)$处形成线性化方程组.\\
在$(\pi,0)$处:
\[
\mathcal{J} = 
\begin{bmatrix}
    0 & 1 \\
    -2\cos\theta & -1
\end{bmatrix}_{(\pi,0)} =
\begin{bmatrix}
    0 & 1 \\
    2 & -1
\end{bmatrix}
\]
\[ \lambda^2 + \lambda - 2 = 0 \]
\[ \lambda = 1 , -2 \]
$\lambda = 1$, $-a_1+a_2 = 0$, 
$\begin{bmatrix}
    1 \\
    1
\end{bmatrix}e^t$ \\
$\lambda = -2$, $2a_1 + a_2 = 0$, 
$\begin{bmatrix}
    1 \\
    -2
\end{bmatrix}e^{-2t}$\\
\begin{figure}[htbp]
    \centering
\includegraphics[width=0.5\textwidth]{C:/Users/37/Desktop/Notes on Common Calculus Equations/Notes and pictures/cwf12.4.3.png}
    \caption{源节点发散}
    \label{fig:cwf12.4.3}
\end{figure}

\newpage
3.构成大图\\
\begin{figure}[htbp]
    \centering
\includegraphics[width=1\textwidth]{C:/Users/37/Desktop/Notes on Common Calculus Equations/Notes and pictures/cwf12.4.4.png}
    \caption{区域性变化}
    \label{fig:cwf12.4.4}
\end{figure}

参考图:
\begin{figure}[htbp]
    \centering
\includegraphics[width=1\textwidth]{C:/Users/37/Desktop/Notes on Common Calculus Equations/Notes and pictures/cwf12.4.5.png}
    \caption{参考图}
    \label{fig:cwf12.4.5}
\end{figure}

至此,图像大致趋势绘制完成.